\chapter{Análisis DAFO del proyecto}
\label{anexo:dafo}

Este anexo presenta un análisis DAFO del proyecto TierraViva, con el objetivo de valorar de forma estructurada sus fortalezas, debilidades, oportunidades y amenazas. El análisis complementa la descripción técnica del prototipo con una visión crítica sobre su viabilidad, sus limitaciones actuales y sus posibles líneas de evolución.

El DAFO se plantea desde una perspectiva técnico-académica, no comercial. Por tanto, las debilidades y amenazas no se interpretan como fallos del proyecto, sino como aspectos reales que deben tenerse en cuenta para evaluar correctamente el alcance del prototipo.

\section{Matriz DAFO}

La Tabla \ref{cuadro:dafo_interno} recoge los factores internos del proyecto: fortalezas y debilidades.

\begin{table}[H]
\begin{center}
\small
\begin{tabular}{|p{0.45\textwidth}|p{0.45\textwidth}|}
\hline
\rowcolor{gray!20}
\textbf{Fortalezas} & \textbf{Debilidades} \\
\hline
\begin{minipage}[t]{0.43\textwidth}
\vspace{0.1cm}
\begin{itemize}
    \item \textbf{F1. Arquitectura modular:} separación clara entre estación de campo, plataforma IoT, almacenamiento local, API y \textit{dashboard}.
    \item \textbf{F2. Decisión local de riego:} la estación puede actuar sobre la bomba sin depender continuamente de la Raspberry Pi.
    \item \textbf{F3. Conectividad móvil:} el uso del módulo LTE A7670E-LASA evita depender de una red WiFi local en la parcela.
    \item \textbf{F4. Prototipo funcional:} integra sensorización, comunicaciones, almacenamiento, visualización y actuación física.
    \item \textbf{F5. Registro histórico:} la Raspberry Pi permite conservar lecturas y facilitar la trazabilidad del sistema.
    \item \textbf{F6. Bajo coste relativo:} se emplean componentes comerciales accesibles y reproducibles.
\end{itemize}
\vspace{0.1cm}
\end{minipage}
&
\begin{minipage}[t]{0.43\textwidth}
\vspace{0.1cm}
\begin{itemize}
    \item \textbf{D1. Dependencia de la señal LTE local:} el funcionamiento remoto requiere señal suficiente, antena adecuada y alimentación estable del módem.
    \item \textbf{D2. Robustez limitada:} el montaje no alcanza todavía el nivel mecánico, eléctrico y ambiental de un producto industrial, ya que el objetivo de esta versión es validar un prototipo académico funcional. Su mejora se plantea como una línea clara de evolución futura.
    \item \textbf{D3. Calibración de sensores:} la humedad del suelo depende del sustrato, la profundidad y la colocación del sensor.
    \item \textbf{D4. Consumo del módulo LTE:} la comunicación móvil puede generar picos de corriente y aumentar el consumo total.
    \item \textbf{D5. Escalabilidad validada parcialmente:} la arquitectura admite más estaciones, pero la validación se centra en un número reducido de nodos.
    \item \textbf{D6. Sin predicción avanzada:} la lógica de riego se basa en reglas y umbrales, no en modelos predictivos.
\end{itemize}
\vspace{0.1cm}
\end{minipage}
\\
\hline
\end{tabular}
\caption{Matriz DAFO del proyecto TierraViva. Factores internos.}
\label{cuadro:dafo_interno}
\end{center}
\end{table}

La Tabla \ref{cuadro:dafo_externo} recoge los factores externos: oportunidades y amenazas.

\begin{table}[H]
\begin{center}
\small
\begin{tabular}{|p{0.45\textwidth}|p{0.45\textwidth}|}
\hline
\rowcolor{gray!20}
\textbf{Oportunidades} & \textbf{Amenazas} \\
\hline
\begin{minipage}[t]{0.43\textwidth}
\vspace{0.1cm}
\begin{itemize}
    \item \textbf{O1. Agricultura de precisión:} el sistema puede ampliarse con más estaciones, sensores y criterios de decisión.
    \item \textbf{O2. Predicción meteorológica:} el riego podría ajustarse usando lluvia prevista, temperatura o evapotranspiración.
    \item \textbf{O3. Optimización del agua:} la automatización basada en humedad real del suelo puede reducir riegos innecesarios.
    \item \textbf{O4. Mejora energética:} el prototipo podría evolucionar hacia alimentación solar, batería y modos de bajo consumo.
    \item \textbf{O5. Sistema de alertas:} se podrían incorporar avisos ante fallos, desconexiones o valores críticos.
    \item \textbf{O6. Dashboard multiestación:} la interfaz puede evolucionar hacia una herramienta de supervisión más completa.
\end{itemize}
\vspace{0.1cm}
\end{minipage}
&
\begin{minipage}[t]{0.43\textwidth}
\vspace{0.1cm}
\begin{itemize}
    \item \textbf{A1. Degradación de comunicación:} pérdidas de señal, incidencias del operador o reinicios del módem pueden retrasar la telemetría.
    \item \textbf{A2. Condiciones ambientales adversas:} lluvia, humedad, polvo o temperatura pueden afectar a sensores, cableado y electrónica.
    \item \textbf{A3. Fallos hidráulicos:} una bomba obstruida, un tubo desconectado o una fuga pueden impedir el riego aunque la lógica funcione.
    \item \textbf{A4. Lecturas erróneas:} un sensor mal colocado o degradado puede provocar decisiones incorrectas.
    \item \textbf{A5. Dependencia de servicios externos:} ThingSpeak introduce dependencia de disponibilidad, límites de uso y cambios de plataforma.
    \item \textbf{A6. Seguridad y credenciales:} una implantación real requeriría reforzar la protección de claves, \textit{endpoints} y canales.
\end{itemize}
\vspace{0.1cm}
\end{minipage}
\\
\hline
\end{tabular}
\caption{Matriz DAFO del proyecto TierraViva. Factores externos.}
\label{cuadro:dafo_externo}
\end{center}
\end{table}

\section{Lectura sintética del análisis}

El análisis muestra que TierraViva tiene una base técnica sólida para un prototipo académico de riego inteligente. Sus principales fortalezas se encuentran en la integración de varias capas: estación embebida, comunicación LTE, plataforma IoT, almacenamiento local, API y \textit{dashboard}.

Las debilidades se relacionan principalmente con la madurez del prototipo: protección ambiental, calibración, consumo energético y validación limitada a un número reducido de nodos. Estas limitaciones no invalidan el proyecto, pero delimitan correctamente su alcance.

Las oportunidades apuntan hacia una evolución natural del sistema: más estaciones, mayor robustez, predicción meteorológica, alertas y optimización energética. Las amenazas se concentran en aspectos típicos de sistemas IoT de campo: comunicación, exposición ambiental, fiabilidad de sensores, dependencia de servicios externos y seguridad.

\section{Estrategias derivadas del DAFO}

Las estrategias se identifican combinando las iniciales de los bloques DAFO: fortalezas y oportunidades (FO), debilidades y oportunidades (DO), fortalezas y amenazas (FA), y debilidades y amenazas (DA).

A partir de la matriz anterior se plantean cuatro estrategias de mejora y consolidación.

\begin{table}[H]
\begin{center}
\small
\begin{tabular}{|p{0.18\textwidth}|p{0.72\textwidth}|}
\hline
\rowcolor{gray!20}
\textbf{Tipo} & \textbf{Aplicación en TierraViva} \\
\hline
\textbf{FO} & Aprovechar la arquitectura modular, la decisión local y el registro histórico para evolucionar hacia un sistema de agricultura de precisión con más estaciones y análisis de tendencias. \\
\hline
\textbf{DO} & Reducir las debilidades mediante calibración de sensores, validación de datos, mejora energética y protección física del montaje. \\
\hline
\textbf{FA} & Usar la decisión local, el apagado seguro y la separación entre supervisión y actuación para mitigar fallos de comunicación, lecturas anómalas o indisponibilidad temporal de servicios externos. \\
\hline
\textbf{DA} & Plantear futuras mejoras centradas en robustez ambiental, seguridad de credenciales, tolerancia a fallos, encapsulado y reducción de dependencias externas. \\
\hline
\end{tabular}
\caption{Estrategias derivadas del análisis DAFO.}
\label{cuadro:estrategias_dafo}
\end{center}
\end{table}

\section{Conclusión}

El análisis DAFO confirma que TierraViva es un proyecto viable, estable, funcional y coherente. Su valor principal reside en integrar hardware, comunicaciones y \textit{software} para resolver un problema real: la supervisión y automatización del riego en un entorno desatendido.

Al mismo tiempo, el DAFO permite reconocer de forma honesta los límites del prototipo. La cobertura móvil local, la calibración de sensores, la exposición ambiental, la seguridad de credenciales y la dependencia de servicios externos son aspectos que deben tratarse con cautela en futuras versiones.